\documentclass[french]{article}

\usepackage[utf8]{inputenc}
\usepackage[T1]{fontenc}
\usepackage{lmodern}
\usepackage{geometry}
\usepackage{graphicx}
\usepackage{babel}
\usepackage{amsmath}
\usepackage{amsthm}
\usepackage{amssymb}
\usepackage{hyperref}
\usepackage{stmaryrd}
\usepackage{enumitem}
\usepackage[most]{tcolorbox}
\usepackage{dsfont}
\usepackage{multirow}
\usepackage{etoc}
\usepackage{titlesec}
\usepackage{esint}
\usepackage{cancel}
\usepackage{mathdots}

\setlist[itemize]{label=$\bullet$}


\renewcommand{\P}{\mathbb{P}}
\newcommand{\N}{\mathbb{N}}
\newcommand{\Z}{\mathbb{Z}}
\newcommand{\Q}{\mathbb{Q}}
\newcommand{\R}{\mathbb{R}}
\newcommand{\C}{\mathbb{C}}
\newcommand{\K}{\mathbb{K}}
\renewcommand{\L}{\mathbb{L}}
\newcommand{\U}{\mathbb{U}}
\newcommand{\st}{^*}
\newcommand{\tq}{\middle|}
\newcommand{\inv}{^{-1}}
\newcommand{\E}{\mathbb{E}}
\newcommand{\F}{\mathbb{F}}
\newcommand{\V}{\mathbb{V}}
\renewcommand{\ker}{\text{Ker}}
\newcommand{\im}{\text{Im}}
\newcommand{\card}{\text{Card}}
\newcommand{\Tr}{\text{Tr}}
\newcommand{\Vect}{\text{Vect}}
\newcommand{\rg}{\text{rg}}
\newcommand{\vertiii}[1]{{\left\vert\kern-0.25ex\left\vert\kern-0.25ex\left\vert #1 \right\vert\kern-0.25ex\right\vert\kern-0.25ex\right\vert}}
\renewcommand{\o}{\text{o}}
\renewcommand{\O}{\text{O}}
\renewcommand{\d}{\text{d}}

\theoremstyle{plain}
\newtheorem*{ind}{Indication}

\theoremstyle{definition}

\newtheorem*{ex}{Exemple}
\newtheorem*{met}{Point méthode}
\newtheorem{exo}{Exercice}
\newtheorem*{rmq}{Remarque}

\theoremstyle{remark}
\newtheorem*{app}{Application}

\newtcolorbox[auto counter]{dfn}{
	enhanced,
	breakable,
	colback=white,
	fonttitle=\sc,
	title={Définition \thetcbcounter}
}

\newtcolorbox[auto counter]{thm}[2][]{enhanced, breakable, colback=white, fonttitle=\sc, title=Théorème~\thetcbcounter~: #2}


\title{Suites et séries numériques}
\date{}

\begin{document}
\tableofcontents
\newpage

\maketitle

Dans tout ce chapitre, $\K$ désigne le corps des réels ou le corps des complexes.

\section{Suites numériques}

\subsection{Rappels}

Revoir dans le cours de première année :
\begin{itemize}
	\item la définition d'une suite convergente, d'une suite tendant vers $+\infty$ ou $-\infty$ ;
	\item les suites monotones ;
	\item les suites adjacentes ;
	\item la comparaison des suites.
\end{itemize}

\begin{ex}
	\begin{enumerate}
		\item $\ln(\cos(1/n))\sim\cos(1/n)-1\sim-1/2n^2$, et pas le contraire.
		\item Équivalent de $\cos(1/n)-1+1/n$ puis de $\cos(1/n)-1+1/n^2$.
		\item $\exp(a_n)\sim\exp(b_n)\iff b_n-a_n\xrightarrow[n\to+\infty]{}0$
	\end{enumerate}
\end{ex}

\begin{thm}{(HP)}%prop
	Si $(a_n)$ et $(b_n)$ sont deux suites à termes strictement positifs équivalentes qui tendent vers $+\infty$ ou qui tendent vers $0$, alors $\ln(a_n)\sim\ln(b_n)$.
\end{thm}

\begin{ex}
	Existence, limite et équivalent de $(x_n)\in[0,1]^\N$ vérifiant $\forall n\in\N,\ x_n=(1-x_n)^n$.
\end{ex}

Exercice Moulin.SSN.05

\subsection{Suites extraites, valeurs d'adhérence}

\begin{dfn}
	\begin{itemize}
		\item On appelle \textbf{extraction} toute application $\varphi$ strictement croissante de $\N$ dans $\N$.
		\item Soit $(u_n)_{n\in\N}$ une suite. On appelle \textbf{suite extraite} ou \textbf{sous-suite} tout suite de la forme $(u_{\varphi(n)})_{n\in\N}$ où $\varphi$ est une extraction.
	\end{itemize}
\end{dfn}

\begin{thm}{}%lem
	Si $\varphi$ est une extraction, alors $\forall n\in\N,\ \varphi(n)\geqslant n$.
\end{thm}

\begin{thm}{}%prop
	Soit $(u_n)_{n\in\N}$ une suite de $\K$ convergeant vers $l\in\K$. Alors toute sous-suite de $u$ converge également vers $l$.
\end{thm}

\begin{rmq}
	De la même façon, si $(u_n)$ est une suite réelle tendant vers $+\infty$, toutes ses sous-suites tendent également vers $=\infty$. Il en est de même pour $-\infty$.
\end{rmq}

\begin{dfn}
	On appelle \textbf{valeur d'adhérence} d'une suite $(a_n)_{n\in\N}$ tout élément de $\K$ qui est limite d'une sous-suite convergente de $(a_n)_{n\in\N}$.
\end{dfn}

\begin{rmq}
	D'après la proposition précédente, une suite convergente ne possède que sa limite comme valeur d'adhérence.
\end{rmq}

\begin{met}
	Pour montrer qu'une suite est divergente, on peut utiliser la contraposée de la remarque précédente : il suffit de montrer que la suite considérée possède au moins deux valeurs d'adhérence.
\end{met}

\begin{ex}
	La suite réelle de terme général $(-1)^n$ est divergente, car elle possède $1$ et $-1$ comme valeurs d'adhérence.
\end{ex}

\begin{exo}
	On considère la suite de terme général $a_n=n(1+(-1)^n)$.
	\begin{enumerate}
		\item Montrer que $0$ est la seule valeur d'adhérence de la suite $(a_n)_{n\in\N}$.
		\item La suite $(a_n)_{n\in\N}$ est-elle convergente ?
	\end{enumerate}
\end{exo}

\paragraph{Attention}
Comme le montre l'exercice précédent, ce n'est pas parce qu'une suite ne possède qu'une seule valeur d'adhérence qu'elle converge !

\paragraph{Rappel}
La convergence des deux suites $(u_{2n})$ et $(u_{2n+1})$ vers la même limite $l$ entraîne la convergence de $(u_n)$ vers $l$.

\begin{exo}
	Soit $\varphi$ et $\psi$ deux applications strictement croissantes de $\N$ dans $\N$. On suppose que les deux sous-suites $(u_{\varphi(n)})$ et $(u_{\psi(n)})$ convergent vers un même élément $l$. Donner une condition nécessaire et suffisante sur $\varphi$ et $\psi$ pour que l'on puisse en déduire que $(u_n)$ converge vers $l$.
\end{exo}

\begin{thm}{Caractérisation ensembliste des valeurs d'adhérence}
	Pour une suite $(u_n)_{n\in\N}$ et un élément $l\in\K$, il est équivalent de dire :
	\begin{enumerate}[label=(\roman*)]
		\item $l$ est une valeur d'adhérence de $(u_n)_{n\in\N}$ ;
		\item $\forall\varepsilon>0,\ \{n\in\N\ :\ \lvert u_n-l\rvert\leqslant\varepsilon\}$ est infini ;
		\item $\forall\varepsilon>0,\ \forall n\in\N,\ \exists p\geqslant n,\ \lvert u_p-l\rvert\leqslant\varepsilon$.
	\end{enumerate}	
\end{thm}

Exercice Moulin.SSN.02

\subsection{Théorème de Bolzano-Weierstrass}

\begin{thm}{Théorème de Bolzano-Weierstrass}%thm
	Toute suite bornée de $\K$ admet une valeur d'adhérence.
\end{thm}

\begin{dfn}{Limites inférieure et supérieure (HP)}
	Soit $u$ une suite réelle.
	\begin{itemize}
		\item Si $u$ est minorée, la suite $(y_n)_{n\in\N}$ définie par : $$y_n=\inf\{u_p\ \mid\ p\geqslant n\}$$ est croissante. Sa limite (éventuellement $+\infty$) est appelée \textbf{limite inférieure} de $u$ et est notée $\underline{\lim}u$. \\
		Si $u$ est non minorée, on pose $\underline{\lim}u=-\infty$.
		\item Si $u$ est majorée, la suite $(z_n)_{n\in\N}$ définie par : $$y_n=\sup\{u_p\ \mid\ p\geqslant n\}$$ est décroissante. Sa limite (éventuellement $-\infty$) est appelée \textbf{limite supérieure} de $u$ et est notée $\overline{\lim}u$.\\
		Si $u$ est non majorée, on pose $\overline{\lim}u=+\infty$.
	\end{itemize}
\end{dfn}

\begin{thm}{(HP)}%prop
	On a $\underline{\lim}u\leqslant\overline{\lim}u$, avec égalité si, et seulement si, $u$ admet une limite.
\end{thm}

\begin{exo}
	Montrer qu'il existe une sous-suite de $u$ qui admet $\underline{\lim}u$ pour limite, et qu'il existe une sous-suite de $u$ qui admet $\overline{\lim}u$ pour limite. \\
	En déduire une autre démonstration du théorème de Bolzano-Weierstrass pour les suites réelles. \\
	Montrer que tout valeur d'adhérence de $u$ est comprise entre $\underline{\lim}u$ et $\overline{\lim}u$.
\end{exo}

\begin{thm}{Critère de Cauchy (HP)}
	Une suite complexe $u$ converge si, et seulement si : $$\forall\varepsilon>0,\ \exists n_0\in\N,\ \forall p\geqslant n_0,\ \forall q\geqslant n_0,\ \lvert u_p-u_q\rvert\leqslant\varepsilon$$
\end{thm}

\section{Séries numériques}

\subsection{Suites et séries}

A toute suite $(u_n)_{n\in\N}\in\K^\N$, on associe la suite $(S_n)_{n\in\N}$ définie par : $$\forall n\in\N,\ S_n=\sum_{k=0}^{n}u_k$$ L'application $u\mapsto S$ est un isomorphisme d'espace vectoriel de $\K^\N$. Sa réciproque est l'application $S\mapsto u$, où $u$ est définie par : $$u_0=S_0\ \ \text{et}\ \ \forall n\in\N\st,\  u_n=S_n-S_{n-1}$$ Pour chaque élément de $\K^\N$, on peut donc considérer au choix la suite $(u_n)_{n\in\N}$ ou la suite $(S_n)_{n\in\N}$ associée.

\begin{dfn}
	Si $u\in\K^\N$, et $S_n=\displaystyle\sum_{p=0}^{n}u_p$, on appelle série de terme général CENSURE. \\
	On la note $\displaystyle\sum u_n$. \\
	Le \textbf{terme général} est donc le scalaire $u_n$ ; les \textbf{sommes partielles} sont les $S_n$.
\end{dfn}

\begin{dfn}
	On dit que la série $\displaystyle\sum u_n$ \textbf{converge} si la suite $(S_n)$ admet une limite dans $\K$, appelée \textbf{somme de la série} et notée $\displaystyle\sum_{n=0}^{+\infty}u_n$. Dans le cas contraire, on dit que la série \textbf{diverge}. On parle de la nature d'une série pour parler de sa convergence ou sa divergence.
\end{dfn}

\paragraph{Notation} On note $S_n(u)$ la $n$-ième somme partielle de la série de terme général $u_n$. Lorsque la série converge, on note aussi $S(u)$ sa somme.

\begin{rmq}
	\begin{itemize}
		\item Lorsque la suite $u$ est définie seulement à partir d'un certain rang $n_0$, la série de terme général $u_n$ est notée $\displaystyle\sum_{n\geqslant n_0}u_n$ ; il s'agit de la suite des sommes partielles définie par : $$\forall n\geqslant n_0,\ S_n=\sum_{p=n_0}^{n}u_p$$
		\item Le caractère asymptotique de la notion de limite permet de voir que pour $n_0\in\N$, une série $\displaystyle\sum u_n$ converge si, et seulement si, la série $\displaystyle\sum_{n\geqslant n_0}u_n$.
	\end{itemize}
\end{rmq}

\paragraph{Suites vs séries}
\begin{itemize}
	\item L'étude d'une série n'est donc autre chose que l'étude d'une suite.
	\item MAIS une série est donnée par son terme général et non par ses sommes partielles, et son étude, spécifique, consiste justement à travailler sur le terme général.
	\item On n'étudie donc pas (sauf rares exceptions) une série par considérations sur la suite de ses sommes partielles. Mieux, pour étudier une suite, on peut considérer une série :
\end{itemize}

\begin{met}
	Pour étudier la convergence d'une suite $(u_n)_{n\in\N}$, il faut penser à étudier éventuellement la convergence de la série $\displaystyle\sum_{n\geqslant 1}(u_{n+1}-u_n)$.
\end{met}

\begin{thm}{}%prop
	Si une série converge, alors son terme général converge vers $0$.
\end{thm}

\paragraph{Attention} La convergence vers $0$ du terme général est une condition nécessaire mais non suffisante de convergence. On pourra par exemple considérer la série dont le terme général est donné par : $u_0=0$ et $\forall n\in\N\st,\ u_n=\sqrt{n}-\sqrt{n-1}$.

\begin{thm}{}%prop
	L'ensemble des suites $(u_n)_{n\in\N}\in\K^\N$ telles que la série $\displaystyle\sum u_n$ converge est un sous-espace vectoriel de $\K^\N$ sur lequel l'application somme est linéaire.
\end{thm}

\begin{thm}{}%cor
	On ne change pas la nature d'une série en ajoutant à son terme général le terme général d'une série convergente.
\end{thm}

\begin{rmq}
	Cette proposition sert pratiquement à "découper" le terme général d'une série en bouts plus simples à étudier, par exemple en faisant un développement asymptotique du terme général.
\end{rmq}

\begin{ex}%exs
	\begin{enumerate}
		\item Une série géométrique $\displaystyle\sum a^n$ converge si, et seulement si, $\lvert a\rvert<1$. Sa somme (si elle commence au rang $n=0$) est alors $1/(1-a)$.
		\item La série $\displaystyle\frac{1}{n}$ diverge : minoration de $S_{2n}-S_n$.
		\item Une série télescopique est une série dont le terme général est naturellement sous la forme $u_{n+1}-u_n$. Ses sommes partielles sont donc les $u_{n+1}-u_n$. Elle converge donc si, et seulement si, la suite $u$ admet une limite finie. Sa somme est alors $\lim u-u_0$.
	\end{enumerate}
\end{ex}

\begin{exo}
	Convergence et somme des séries :
	$$\sum\frac{1}{n(n+1)}\ \ \text{et}\ \ \sum\arctan\left(\frac{1}{1+n+n^2}\right)$$
\end{exo}

\begin{dfn}
	Lorsque la série $\displaystyle\sum u_n$, on note $\displaystyle R_n(u)=\sum_{p=n+1}^{+\infty}u_p$, qu'on appelle \textbf{reste} au rang $p$.
\end{dfn}

\begin{rmq}
	Dans ce cas, la suite $(R_n(u))$ converge vers $0$ comme le prouve la relation : $$R_n(u)=S(u)-S_n(u)$$
\end{rmq}

\begin{exo}
	Que penser de l'affirmation "la suite $(R_n(u))$ converge vers $0$ puisque $\displaystyle\sum u_n$ converge" ?
\end{exo}

\subsection{Séries à terme général positif}

\begin{thm}{}%prop
	Soit $\displaystyle\sum u_n$ une série à terme général réel positif ou nul. Alors les sommes partielles convergent ou tendent vers $+\infty$.
\end{thm}

\begin{thm}{Comparaison de séries à termes généraux positifs}
	Soit $\displaystyle\sum u_n$ et $\displaystyle\sum v_n$ deux séries vérifiant : $$\exists n_0\in\N,\ \forall n\geqslant n_0,\ 0\leqslant u_n\leqslant v_n$$
	\begin{itemize}
		\item Si $\displaystyle\sum v_n$ converge, alors $\displaystyle\sum u_n$ converge.
		\item Si $\displaystyle\sum u_n$ diverge, alors $\displaystyle\sum v_n$ diverge.
	\end{itemize}
\end{thm}

\subsection{Convergence absolue}

\begin{dfn}
	On dit que la série $\displaystyle\sum u_n$ est \textbf{absolument convergente}, ou \textbf{converge absolument}, si la série réelle positive $\displaystyle\sum \lvert u_n\rvert$ est convergente. Cela équivaut à la sommabilité de $(u_n)_{n\in\N}$.
\end{dfn}

\begin{thm}{}
	Toute série réelle ou complexe absolument convergente est convergente.
\end{thm}

\begin{dfn}
	Une série qui converge mais qui n'est pas absolument convergente est dite \textbf{semi-convergente}.
\end{dfn}

\begin{ex}
	La suite $(-1)^n/n$ converge vers $0$. C'est la suite des sommes partielles dont le terme général $u_n$ vérifie $\lvert u_n\rvert\geqslant1/n$ donc qui n'est pas absolument convergente.
\end{ex}

\section{Étude de convergence de séries}

\subsection{Théorèmes de comparaison}

\begin{thm}{}%prop
	Soit $(u_n)$ et $(v_n)$ deux séries telles que $u_n=\O(v_n)$. Si $(v_n)$ est sommable, alors $(u_n)$ est sommable.
\end{thm}

\begin{rmq}
	Évidemment, le résultat subsiste si l'on a $u_n=\o(v_n)$.
\end{rmq}

\begin{ex}
	Soit $\displaystyle\sum u_n$ et $\displaystyle\sum v_n$ deux séries à terme général réel. Si $u_n\sim v_n$ et si $(v_n)$ est positive, alors $\displaystyle\sum u_n$ et $\displaystyle\sum v_n$ sont de même nature.
\end{ex}

Exercice Moulin.SSN.22

\subsection{Séries de référence}

\begin{rmq}
	Pour démontrer la convergence (ou la divergence) d'une série à terme général positif à l'aide des théorèmes de comparaison ci-dessus, on utilise des séries de référence : séries géométriques, ou séries de Riemann.
\end{rmq}

\begin{thm}{Séries de Riemann}%prop
	La suite $\displaystyle\left(\frac{1}{n^\alpha}\right)_{n\geqslant 1}$ est sommable si, et seulement si, $\alpha >1$.
\end{thm}
\begin{proof}
	Trouver un équivalent $(n+1)^\beta-n^\beta$.
\end{proof}

\begin{met}
	On a les différents cas.
	\begin{itemize}
		\item Si $\exists\alpha >1 \ : \ n^\alpha u_n\to 0$, alors $\sum u_n$ converge absolument.
		\item Si $\exists\alpha \leqslant1 \ : \ n^\alpha u_n\to+\infty$, alors $\sum u_n$ diverge.
		\item Si $\exists M\geqslant m>0\ : \ m\leqslant n^\alpha u_n\leqslant M$, alors $\sum u_n$ est de même nature que $\sum n^{-\alpha}$.
	\end{itemize}
\end{met}

\begin{ex}
	$\displaystyle\sum\frac{\ln(n)^\beta}{n^\alpha}$, $\displaystyle\sum e^{-n^\alpha}$ et $\displaystyle\sum\left(\cos\left(\frac{1}{n}\right)\right)^{n^\alpha}$
\end{ex}

\subsection{Règle de d'Alembert}

\begin{thm}{Comparaison logarithmique}%prop
	Si, à partir d'un certain rang, $u_n$ et $v_n$ sont strictement positifs et si $\displaystyle\frac{u_{n+1}}{u_n}\leqslant\frac{v_{n+1}}{v_n}$, alors $u_n=\O(v_n)$, et donc :
	\begin{itemize}
		\item la convergence de $\displaystyle\sum v_n$ entraîne la convergence de $\displaystyle\sum u_n$ ;
		\item la divergence de $\displaystyle\sum u_n$ entraîne la divergence.
	\end{itemize}
\end{thm}

\begin{thm}{}%cor
	Soit $(u_n)$ une suite à valeurs strictement positives. Si à partir d'un certain rang :
	\begin{enumerate}
		\item $\displaystyle\frac{u_{n+1}}{u_n}\leqslant k<1$, alors $\displaystyle\sum u_n$ converge.
		\item $\displaystyle\frac{u_{n+1}}{u_n}\geqslant 1$, alors $\displaystyle\sum u_n$ diverge.
	\end{enumerate}
\end{thm}

\begin{thm}{}%cor
	Soit $(u_n)$ une suite à valeurs strictement positives. Si :
	\begin{enumerate}
		\item $\displaystyle\underset{+\infty}{\lim}\frac{u_{n+1}}{u_n}<1$, alors $\displaystyle\sum u_n$ converge.
		\item $\displaystyle\underset{+\infty}{\lim}\frac{u_{n+1}}{u_n}>1$, alors $\displaystyle\sum u_n$ diverge.
	\end{enumerate}
\end{thm}

\begin{ex}%exs
	$\displaystyle\sum\frac{x^n}{n!}$, $\displaystyle\sum\binom{2n}{n}x^n$
\end{ex}

Exercice Moulin.SSN.11

\subsection{Règle de Raabe-Duhamel (HP)}

\paragraph{Attention} Si $\displaystyle\frac{u_{n+1}}{u_n}\to 1$, on ne peut rien conclure directement. La règle suivante généralise la règle de d'Alembert par comparaison logarithmique avec des séries de Riemann.

\begin{thm}{Règle de Raabe-Duhamel (HP)}%prop
	Soit $(u_n)_{n\in\N}$ une suite à valeurs strictement positives vérifiant :
	$$\frac{u_{n+1}}{u_n}=1-\frac{\alpha}{n}+\o\left(\frac{1}{n}\right)$$
	Alors :
	\begin{itemize}
		\item si $\alpha>1$, alors $\displaystyle\sum u_n$ converge.
		\item si $\alpha<1$, alors $\displaystyle\sum u_n$ diverge.
	\end{itemize}
\end{thm}

Lorsque $\alpha=1$, on ne peut rien conclure. Les séries de Bertrand $\displaystyle\sum u_n$ avec : $$u_n=\frac{1}{n\ln(n)^\alpha}$$ vérifient : $\displaystyle\frac{u_{n+1}}{u_n}=1-\frac{1}{n}+\frac{\alpha}{n\ln(n)}+\o\left(\frac{1}{n\ln(n)}\right)$ alors que, suivant les valeurs de $\alpha$, la série $\displaystyle\sum u_n$ peut converger ou diverger. \\

Néanmoins : 
\begin{ex}
	\begin{enumerate}
		\item Soit $(u_n)_{n\in\N}$ une suite à valeurs strictement positives vérifiant : $$\frac{u_{n+1}}{u_n}=1+\frac{\alpha}{n}+\O\left(\frac{1}{n^2}\right)$$ Alors il existe $K>0$ telle que $u_n\sim Kn^\alpha$.
		\item Cas $\alpha=1$ : si $\displaystyle\frac{u_{n+1}}{u_n}=1-\frac{1}{n}+\O\left(\frac{1}{n^2}\right)$, alors $\displaystyle\sum u_n$ diverge.
		\item Formule de Stirling.
	\end{enumerate}
\end{ex}

\begin{exo}
	Soit $I_n=\displaystyle\int_{0}^{\pi/2}\cos^n(t)dt$.
	\begin{enumerate}
		\item Montrer que la suite $(I_n)$ est croissante.
		\item Trouver une relation de récurrence entre $I_n$ et $I_{n+2}$.
		\item En déduire $(n+1)I_nI_{n+1}=\pi/2$.
		\item Donner un équivalent de $I_n$.
		\item Donner une expression de $I_{2n}$ ou $I_{2n+1}$ à l'aide de factorielles et en déduire la formule de Stirling.
	\end{enumerate}
\end{exo}

\subsection{Technique de comparaison série-intégrale}

La technique de comparaison série-intégrale a déjà été vue en première année. Nous donnons ici des approfondissements.

\subsubsection*{Utilisation de la monotonie}

\begin{thm}{}
	Soit $f\in\mathcal{CM}([0,+\infty[,\R_+)$ une fonction décroissante et intégrable. \\
	Alors la série $\displaystyle\sum f(n)$ est convergente et, en posant $I=\displaystyle\int_{0}^{+\infty}f(t)dt$, on a : $$I\leqslant\sum_{n=0}^{+\infty}f(n)\leqslant I+f(0)\ \ \text{et donc}\ \ \left\lvert I-\sum_{n=0}^{+\infty}f(n)\right\rvert\leqslant f(0)$$
\end{thm}
\begin{proof}
	Faire une comparaison série-intégrale pour les sommes partielles et passer à la limite.
\end{proof}

\begin{rmq}
	On peut sans difficulté adapter le résultat précédent pour encadrer les restes de certaines séries convergentes. En effet, si $f\in\mathcal{CM}([n_0,+\infty[,\R_+)$ est une fonction décroissante et intégrable, alors la série $\displaystyle\sum f(n)$ est convergente et :
	$$\forall n\geqslant n_0,\ \int_{n}^{+\infty}f(t)dt\leqslant\sum_{k=n}^{+\infty}f(k)\leqslant f(n)+ \int_{n}^{+\infty}f(t)dt$$
\end{rmq}

\subsubsection*{Interprétation géométrique}

Pour tout $n\in\N$, la quantité $\displaystyle\left\lvert f(n)-\int_{n}^{n+1}f(t)dt\right\rvert$ représente l'aire géométrique de la portion de graphe située entre la droite d'équation $y=f(n)$ et la courbe d'équation $y=f(x)$ pour des abscisses $x$ décrivant $[n,n+1]$. \\

La quantité positive $\displaystyle\left\lvert\int_{0}^{+\infty}f(t)dt-\sum_{n=0}^{+\infty}f(n)\right\rvert$ est la somme de ces aires, que l'on majore par $f(0)$, aire d'un rectangle de largeur $1$ et de hauteur $f(0)$ dans lequel on peut déplacer les hypographes qui correspondent aux aires précédentes.

\begin{ex}
	$\displaystyle\sum_{k=n}^{+\infty}\frac{1}{k^\alpha}\sim\frac{1}{(\alpha-1)n^{\alpha-1}}$
\end{ex}

La propriété suivante traite le cas d'une fonction continue par morceaux, positive, décroissante et intégrable sur l'intervalle ouvert $]0,+\infty[$. On se donne donc une telle fonction $f$.

\begin{thm}{}%prop
	La série $\displaystyle\sum f(n)$ est convergente et l'on a : $$\left\lvert\int_{0}^{+\infty}f(t)dt-\sum_{n=1}^{+\infty}f(n)\right\rvert\leqslant\int_{0}^{1}f(t)dt$$
\end{thm}

Pour tout $n\in\N\st$, la quantité $\displaystyle\left\lvert f(n)-\int_{n-1}^{n}f(t)dt\right\rvert$ représente l'aire géométrique de la portion de graphe située entre la droite d'équation $y=f(n)$ et la courbe d'équation $y=f(x)$ pour des abscisses $x$ décrivant $[n-1,n]$. \\

La quantité positive $\displaystyle\left\lvert\int_{0}^{+\infty}f(t)dt-\sum_{n=0}^{+\infty}f(n)\right\rvert$ est la somme de ces aires, que l'on majore par $\displaystyle\int_{0}^{1}f(t)dt$.

\subsubsection*{Approfondissement : comparaison série-intégrale sans monotonie}

Dans l'exemple qui suit, on traite le cas d'une fonction positive croissante puis décroissante. Cette situation peut se généraliser.

\begin{ex}
	Pour $x>0$, on pose $u_n(x)=\sqrt{n}e^{-x\sqrt{n}}$. Équivalent en $0^+$ de $f(x)=\displaystyle\sum_{n=0}^{+\infty}u_n(x)$.
\end{ex}

\begin{exo}
	Pour tout $n\geqslant 2$, on pose $u_n=\displaystyle\int_{k=1}^{n-1}\frac{1}{\sqrt{k(n-k)}}$. \\
	En effectuant une comparaison avec l'intégrale $\displaystyle\int_{0}^{n}\frac{dt}{\sqrt{t(n-t)}}$, montrer que la suite $(u_n)$ est convergente et préciser sa limite.
\end{exo}

\section{Sommation des relations de comparaison}

\subsection{Cas convergent}

\begin{thm}{}
	Soit $\displaystyle\sum u_n$ et $\displaystyle\sum v_n$ deux séries numériques, la série $\displaystyle \sum v_n$ étant convergente et à terme général \textbf{réel positif}.
	\begin{itemize}
		\item Si $u_n=\O(v_n)$, alors $R_n(u)=\O(R_n(v))$.
		\item Si $u_n=\o(v_n)$, alors $R_n(u)=\o(R_n(v))$.
		\item Si $u_n\sim v_n$, alors $R_n(u)\sim R_n(v)$.
	\end{itemize}
\end{thm}

\begin{ex}
	Équivalent de $\displaystyle\sum_{k=n}^{+\infty}\frac{1}{k^2}$.
\end{ex}

\paragraph{Attention}
Ne pas oublier de vérifier le caractère positif (au moins à partir d'un certain rang) de la série de référence $\displaystyle\sum v_n$.

\begin{ex}
	Équivalent puis développement asymptotique à deux termes de $\displaystyle\sum_{k=n}^{+\infty}\frac{\ln(k+1)-\ln(k)}{k}$.
\end{ex}

\begin{met}
	Lorsque $u_{n+1}/u_n$ admet une limite non nulle différente de $1$, pour trouver un équivalent du reste, on peut sommer l'équivalent $u_{n+1}\sim lu_n$. Si $(u_n)$ et positive sommable et vérifie ces conditions, on peut montrer alors que $$R_n\sim\frac{u_n}{1-l}$$
\end{met}

\begin{exo}
	Équivalent, développement asymptotique à $2$ termes, puis développement asymptotique à $2026$ termes de $\displaystyle\sum_{k=n}^{+\infty}\frac{1}{k!}$.
\end{exo}

"Chaque année c'est plus difficile !"

\subsection{Cas divergent}

\begin{thm}{}
	Soit $\displaystyle\sum u_n$ et $\displaystyle\sum v_n$ deux séries numériques, la série $\displaystyle \sum v_n$ étant divergente et à terme général \textbf{réel positif}.
	\begin{itemize}
		\item Si $u_n=\O(v_n)$, alors $S_n(u)=\O(S_n(v))$.
		\item Si $u_n=\o(v_n)$, alors $S_n(u)=\o(S_n(v))$.
		\item Si $u_n\sim v_n$, alors $S_n(u)\sim S_n(v)$.
	\end{itemize}
\end{thm}

\begin{ex}
	\begin{enumerate}
		\item Soit $(u_n)_{n\in\N}$ une suite telle que $u_n\sim\displaystyle\frac{1}{n^\alpha}$ avec $\alpha<1$. Alors : $S_n(u)\sim\displaystyle\frac{1}{1-\alpha}n^{1-\alpha}$.
		\item Développement asymptotique de $H_n=\displaystyle\sum_{k=1}^{n}\frac{1}{k}$ (utiliser plutôt $H_{n-1}$). \\
		On trouve : $H_n=\displaystyle\ln(n)+\gamma+\frac{1}{2n}-\frac{1}{12n^2}+\O\left(\frac{1}{n^3}\right)$.
	\end{enumerate}
\end{ex}

\paragraph{Attention}
Ne pas oublier de vérifier le caractère positif (au moins à partir d'un certain rang) de la série de référence $\displaystyle\sum v_n$.

\begin{thm}{Théorème de Cesàro - cas fini}%prop
	Soit $u\in\K^\N$ une suite convergente, de limite $l$. Alors on a : 
	$$\frac{1}{n+1}\sum_{k=0}^{n}u_k\xrightarrow[n\to+\infty]{}l$$
\end{thm}

\begin{thm}{Théorème de Cesàro - cas infini}%prop
	Soit $u\in\K^\N$ une suite convergente, de limite $+\infty$. Alors on a : 
	$$\frac{1}{n+1}\sum_{k=0}^{n}u_k\xrightarrow[n\to+\infty]{}+\infty$$
\end{thm}

\begin{rmq}
	En remplaçant $u_n$ par $-u_n$, on obtient un résultat similaire pour une suite qui diverge vers $-\infty$.
\end{rmq}

\subsection{Application à l'étude asymptotique de suites récurrentes}

\begin{ex}
	Convergence et développement asymptotique d'une suite $(u_n)_{n\in\N}$ vérifiant : $$\forall n\in\N,\ u_{n+1}=\sin(u_n)$$
\end{ex}

\begin{ex}
	Même étude avec $u_{n+1}=\displaystyle\frac{1}{2}\arctan(u_n)$.
\end{ex}

\section{Étude de convergence de séries à termes quelconques}

Lorsque la série n'est pas à terme général réel positif ou qu'elle ne converge pas absolument, les résultats précédents ne peuvent pas être utilisées. On dispose néanmoins des techniques suivantes pour en étudier la convergence. \\
Hormis le théorème des séries alternées, aucun des résultats ci-dessous n'est au programme. En revanche, les techniques sont à connaître et à savoir mettre en œuvre pour résoudre les exercices.

\subsection{Séries alternées}

\begin{dfn}
	On dit que la série réelle $\displaystyle\sum u_n$ est \textbf{alternée} si $(-1)^nu_n$ est un réel de signe constant (au sens large).
\end{dfn}

\begin{thm}{Théorème des séries alternées}%thm
	Une série alternée dont le terme général décroît en module et tend vers $0$ est convergente. De plus, $R_n(u)$ est du signe de son premier terme $u_{n+1}$, et l'on a: $$\forall n\in\N,\ \lvert R_n(u)\rvert\leqslant\lvert u_{n+1}\rvert$$
\end{thm}

\paragraph{Attention} Si le terme général ne décroît pas en module, la série peut très bien diverger. par exemple pour : $$u_{2n}=\frac{1}{2^n}\ \ \text{et}\ \ u_{2n+1}=-\frac{1}{2n+1}$$

\begin{ex}
	Sous les mêmes hypothèses, les sommes partielles sont du signe du premier terme.
\end{ex}

\subsection{Utilisation d'un développement asymptotique du terme général (HP)}

\begin{rmq}
	Si le terme général d'une série alternée tend vers $0$ mais ne décroît pas en module, on peut faire un développement limité du terme général.
\end{rmq}

\begin{ex}
	Comportement de $u_n=\displaystyle\frac{(-1)^n}{n^\alpha+(-1)^n}$ en fonction de $\alpha\in\R$.
\end{ex}

\subsection{Sommation par tranches (HP)}

\begin{ex}
	$u_n=\displaystyle\frac{(-1)^n}{n+(-1)^n}$ en regroupant par tranches de $2$.
\end{ex}

Soit $\varphi$ une extraction telle que $\varphi(0)=0$ et $\displaystyle\sum u_n$ une série. On définit : $$v_n=\sum_{p=\varphi(n)}^{\varphi(n+1)-1}u_p$$
La série $\displaystyle\sum v_n$ s'appelle série regroupée de la série $\displaystyle\sum u_n$. On a $S_n(v)=S_{\varphi(n+1)-1}(u)$. \\
Les sommes partielles de la série regroupée forment donc une sous-suite de la suite des sommes partielles de la série de départ. La convergence de la série $\displaystyle\sum u_n$ entraîne ainsi la convergence de $\displaystyle\sum c_n$, avec égalité des sommes, mais al réciproque est évidemment fausse puisque cela signifierait que la convergence d'une sous-suite suffirait à prouver la convergence de la suite.

\begin{ex}
	$u_n=(-1)^n$
\end{ex}

La réciproque devient vraie dans les deux cas suivants.

\begin{thm}{}%prop
	La convergence d'une série regroupée à terme général \textbf{positif} entraîne la convergence de la série.
\end{thm}

Dans le cas général, on a besoin d'une hypothèse supplémentaire.

\begin{thm}{}%prop
	Avec les notations précédentes, si $\sum_{p=\varphi(n)}^{\varphi(n+1)-1}\lvert u_p\rvert$ tend vers $0$ quand $n$ tend vers $+\infty$, alors la convergence de $\displaystyle\sum v_n$ entraîne celle de $\displaystyle\sum v_n$ et on a : $\displaystyle\sum_{n=0}^{+\infty}u_n=\sum_{n=0}^{+\infty}v_n$.
\end{thm}

Cette hypothèse est réalisée en particulier dans chacun des cas suivants :
\begin{itemize}
	\item $u_n\to 0$ et $\varphi(n+1)-\varphi(n)$ est bornée ;
	\item la suite est réelle et le termes de chaque tranche sont du même signe.
\end{itemize}

\begin{ex}
	$\displaystyle\sum\frac{e^{2in\pi/3}}{n}$
\end{ex}

Exercice Moulin.SSN.32 (on pourra remarquer que $1=-1/2+3/2$)

\subsection{Transformation d'Abel (HP)}

La transformation d'Abel est l'équivalent discret de l'intégration par parties et, de même que celle-ci permet de montrer une semi-convergence d'intégrales, celle-là est utile pour montrer une semi-convergence de séries.

\begin{itemize}
	\item La \textbf{dérivée} d'une suite $(u_n)_{n\in\N}$ est la suite $(u_{n+1}-u_n)_{n\in\N}$, ou alternativement $(u_{n}-u_{n-1})_{n\in\N}$ avec la convention $u_{-1}=0$.
	\item Une \textbf{primitive} d'une suite $(u_n)_{n\in\N}$ est une suite dont la dérivée est $(u_n)_{n\in\N}$, par exemple la suite des sommes partielles de $\displaystyle\sum u_n$ ou, en cas de convergence, l'opposée de la suite des restes.
\end{itemize}

\begin{exo}
	Soit $(a_n)$ et $(b_n)$ deux suites à valeurs dans $\K$ et $(A_n)$ une suite vérifiant $\forall n\in\N,\ a_{n}=A_n-A_{n-1}$ (en définissant $A_{-1}$ pour que la relation soit vérifiée pour $n=0$). Montrer : $$\sum_{k=0}^{n}a_kb_k=A_nb_n-A_{-1}b_0-\sum_{k=0}^{n-1}A_k(b_{k+1}-b_{k})$$
\end{exo}

\begin{thm}{}
	Si les sommes partielles de $\displaystyle\sum a_n$ sont bornées et si la série $\displaystyle\sum (b_{n+1}-b_n)$ converge absolument (par exemple si $(b_n)$ tend vers $0$ en décroissant), alors la série $\displaystyle\sum a_nb_n$ converge.
\end{thm}

\begin{exo}
	Donner une majoration du module du reste.
\end{exo}

\begin{ex}
	Série de terme général $\displaystyle\frac{e^{in\theta}}{n^\alpha}$ avec $\theta\ne 0[2\pi]$ et $0<\alpha\leqslant 1$.
\end{ex}

\subsection{Utilisation d'une primitive (HP)}

Pour étudier une série $\displaystyle\sum f(n)$ semi-convergente, on peut considérer une primitive $F$ de $f$ et montrer que la série de terme général $f(n)-(F(n+1)-F(n))$ converge (absolument). La série $\displaystyle\sum f(n)$ converge alors si, et seulement si, $F(n)$ admet une limite finie.

\begin{ex}
	$\displaystyle\sum n^z$ pour $z\in\C$.
\end{ex}

Lorsque $f$ est un produit, on peut calculer une primitive de $f$ à l'aide d'une intégration par parties (dans le bons sens !). On conserve le terme tout intégré en s'assurant que le deuxième est bien négligeable devant celui-ci.

\begin{ex}
	$\displaystyle\sum\frac{e^{i\sqrt{n}}}{n}$
\end{ex}

Exercice Moulin.SSN.36

\end{document}